% Chunk: Chapter 12 introduction.
% Source: Weinberg QFT Vol. I, physical p. 522 (printed p. 499).
% Coverage: chapter-opening prose; citation [1]; no equations, figures, tables, or footnotes.
% Uncertainty: none.

We saw in the previous chapter that calculations in quantum electrodynamics
involving one-loop graphs yield divergent integrals over momentum space, but
that these infinities cancel when we express all parameters of the theory in
terms of ``renormalized'' quantities, such as the masses and charges that are
actually measured. In 1949 Dyson~\hyperref[ref:12.1]{[1]} sketched a proof that
this cancellation would take place to all orders in quantum electrodynamics.
It was immediately apparent (and will be shown here in Sections 12.1 and 12.2)
that Dyson's arguments apply to a larger class of theories with finite numbers
of relatively simple interactions, the so-called renormalizable theories, of
which quantum electrodynamics is just one simple example.

For some years it was widely thought that any sensible physical theory would
have to take the form of a renormalizable quantum field theory. The requirement
of renormalizability played a crucial role in the development of the modern
``standard model'' of weak, electromagnetic, and strong interactions. But as
we shall see here, the cancellation of ultraviolet divergences does not really
depend on renormalizability; as long as we include every one of the infinite
number of interactions allowed by symmetries, the so-called non-renormalizable
theories are actually just as renormalizable as renormalizable theories.

It is generally believed today that the realistic theories that we use to
describe physics at accessible energies are what are known as ``effective
field theories.'' As discussed in Section 12.3, these are low-energy
approximations to a more fundamental theory that may not be a field theory at
all. Any effective field theory necessarily includes an infinite number of
non-renormalizable interactions. Nevertheless, as discussed in Sections 12.3
and 12.4, we expect that at sufficiently low energy all the non-renormalizable
interactions in such effective field theories are highly suppressed.
Renormalizable theories like quantum electrodynamics and the standard model
thus retain their special status in physics, though for reasons that are
somewhat different from those that originally motivated the assumption of
renormalizability in these theories.
